% ECONOMICS_PROBLEM: {"id": "C65-37", "title": "Uniqueness for an oscillatory generator of linear growth", "jel_code": "C65", "source_id": "OP-0564"}
% FIRST_POSTED: 2026-10-09
% ATTEMPT_STATUS: unresolved
% ENTRY_KIND: research_attempt
% !TeX program = pdflatex
% Self-contained: no external input or bibliography files required.
\documentclass[11pt,a4paper]{article}
\usepackage[T1]{fontenc}
\usepackage[utf8]{inputenc}
\usepackage{lmodern}
\usepackage[margin=24mm]{geometry}
\usepackage{amsmath,amssymb,amsthm,mathtools}
\usepackage{booktabs,longtable,array,enumitem,microtype,xurl,hyperref}
\hypersetup{colorlinks=true,linkcolor=blue,citecolor=blue,urlcolor=blue,pdftitle={Checked research attempts: nine economics and stochastic-analysis problems}}
\setlength{\parindent}{0pt}
\setlength{\parskip}{3pt}
\setlength{\emergencystretch}{3em}
\widowpenalty=10000\clubpenalty=10000\displaywidowpenalty=10000
\setlist[enumerate]{leftmargin=*,itemsep=2pt,topsep=3pt}
\newcommand{\R}{\mathbb R}
\newcommand{\N}{\mathbb N}
\newcommand{\E}{\mathbb E}
\newcommand{\Pp}{\mathbb P}
\newcommand{\Q}{\mathbb Q}
\newcommand{\F}{\mathcal F}
\newcommand{\ind}{\mathbf 1}
\newcommand{\cE}{\mathcal E}
\newcommand{\Law}{\operatorname{Law}}
\newcommand{\sgn}{\operatorname{sgn}}
\newcommand{\Var}{\operatorname{Var}}
\newcommand{\dist}{\operatorname{dist}}
\newcommand{\KL}{\operatorname{KL}}
\newcommand{\TV}{\operatorname{TV}}
\newcommand{\Lip}{\operatorname{Lip}}
\DeclareMathOperator*{\esssup}{ess\,sup}
\DeclareMathOperator*{\argmax}{arg\,max}
\newtheorem{theorem}{Theorem}
\newtheorem{lemma}[theorem]{Lemma}
\newtheorem{proposition}[theorem]{Proposition}
\newtheorem{corollary}[theorem]{Corollary}
\theoremstyle{definition}
\newtheorem{example}[theorem]{Example}
\newcommand{\report}[3]{\clearpage\setcounter{theorem}{0}\renewcommand{\thetheorem}{#1.\arabic{theorem}}\renewcommand{\theHtheorem}{#1.\arabic{theorem}}\section*{#1: #2}\addcontentsline{toc}{section}{#1: #2}\textbf{Outcome:} #3.\par\textbf{Tools:} web browsing available; Python computation available. Literature checked on 9 October 2026.}
\newcommand{\stage}[2]{\subsection*{#1) #2}}
\newcommand{\unresolved}{\textbf{LABEL: UNRESOLVED.}}
\newcommand{\disproof}{\textbf{LABEL: FULL SOLUTION. SUBLABEL: COUNTEREXAMPLE/DISPROOF.}}

\begin{document}
\section*{C65-37: Uniqueness for the oscillatory generator \texorpdfstring{$|z|\sin|z|$}{|z| sin |z|}}
\textbf{Outcome:} UNRESOLVED; bounded-terminal and bounded-reference uniqueness proved.\par\textbf{Tools:} web browsing available; Python computation available. Literature checked on 9 October 2026.
\subsection*{1) FORMAL RESTATEMENT}
\paragraph{Brownian conventions used in this report.}
Fix $T>0$ and an integer $d\ge1$. The probability space carries a $d$-dimensional Brownian motion $B$ with its completed, right-continuous natural filtration; $\F=\F_T$ and $\F_0$ is trivial modulo null sets. A scalar solution consists of a continuous adapted real process $Y$ and a predictable $\R^d$-valued process $Z$ such that
\[
 \int_0^T |Z_t|^2dt+\int_0^T|g(t,Y_t,Z_t)|dt<\infty\quad\Pp\text{-a.s.},
\]
and the displayed BSDE holds simultaneously for all times outside one null set. Equality of solutions means indistinguishability of $Y$ and $dt\otimes\Pp$-a.e. equality of $Z$. A process $X$ is of class (D) when $\{X_\tau:\tau\le T\text{ a stopping time}\}$ is uniformly integrable. Logs are natural, norms Euclidean, and $\sgn(0)=0$. At $T=0$ the assertions reduce to $Y_0=\xi$, with no nontrivial $Z$ component.

\paragraph{Precisely stated probabilistic tools.}
We use It\^o's formula and its Tanaka version for continuous semimartingales, Brownian martingale representation, and the following forms of change of measure. If $\theta$ is predictable and $\E\exp(\frac12\int_0^T|\theta|^2)<\infty$, Novikov's theorem says that
$D=\cE(\int\theta\,dB)$ is a uniformly integrable martingale; under $d\Q=D_Td\Pp$, $B-\int\theta dt$ is Brownian. Here $\cE(M)=\exp(M-\langle M\rangle/2)$. Also, for a continuous Brownian martingale, the BMO hypothesis
\[
 \esssup_{\tau\le T}\E\left[\left.\int_\tau^T|\theta_t|^2dt\right|\F_\tau\right]<\infty
\]
implies that $D$ is uniformly integrable and $D_T\in L^q$ for some $q>1$; this is the continuous-martingale BMO exponential theorem \cite{Kazamaki}. Each use below verifies the relevant integrability hypothesis. A local submartingale or supermartingale of class (D) is a true submartingale or supermartingale: apply the localized conditional inequalities and pass to the limit using uniform integrability. A local martingale of class (D) therefore equals the conditional expectation of its terminal value.

Fix $p>1$. For every $\xi\in L^p(\F_T)$, must any two solutions of
\[
 Y_t=\xi+\int_t^T|Z_s|\sin|Z_s|\,ds-\int_t^TZ_s\,dB_s
\]
coincide whenever both $|Y|^p$ processes are of class (D)? The generator is exactly $g(z)=|z|\sin|z|$, not a globally Lipschitz approximation. Class (D) is imposed on the $p$th power; merely assuming that $Y$ is of class (D) would change the statement. Equality of $Z$ is only $dt\otimes\Pp$-a.e., as specified in the Brownian conventions.

Stress points are unbounded oscillatory slopes, uncontrolled difference densities, and the fixed moment order $p$.

\subsection*{2) QUICK LITERATURE/CONTEXT CHECK}
This is Problem 5.3 in Fan--Hu--Tang \cite{FHT}. Its power class-(D) condition matches the attachment. The comparison theorem referenced there requires a uniform continuity property that this generator lacks. Our bounded-data and bounded-reference arguments below verify their own change-of-measure hypotheses rather than invoking that inapplicable comparison theorem.

\subsection*{3) ATTACK PLAN}
\textbf{Proof routes.} Linearize differences and seek a density strong enough for the given $p$; use the linear-growth bound to obtain a bounded-data estimate; compare an arbitrary solution to one with bounded $Z$.

\textbf{Disproof routes.} Test zero and affine terminal values; search pairs of integrands near successive oscillations; test whether localization can preserve terminal coincidence and the power class-(D) requirement simultaneously. The first route is ruled out by the partial theorems; no valid pair for the others is constructed.

\textbf{Landscape and tools.} Radial calculus bounds slopes; oscillatory sequences test uniform continuity; Girsanov removes drifts; H\"older transfers stopping integrability; conditional energy estimates yield BMO; BMO exponentials justify sampling; localization exposes the remaining tail gap.

\subsection*{4) WORK}
\begin{lemma}[Growth and nonuniform continuity]
For all $z,w$, $|g(z)|\le|z|$ and
\[
 |g(z)-g(w)|\le(1+\max\{|z|,|w|\})|z-w|.
\]
The function is not uniformly continuous on $\R^d$.
\end{lemma}
\begin{proof}
For $f(r)=r\sin r$, $|f'(r)|=|\sin r+r\cos r|\le1+r$. Apply the mean value theorem between $|z|$ and $|w|$ and then the reverse triangle inequality. For nonuniform continuity, take points on a fixed positive unit ray with radii
$r_k=2\pi k$ and $s_k=r_k+r_k^{-1}$. Their distances tend to zero, but
\[
 f(s_k)-f(r_k)=(r_k+r_k^{-1})\sin(r_k^{-1})\longrightarrow1,
\]
using $\sin t/t\to1$. This contradicts the defining uniform-continuity implication.
\end{proof}

\begin{lemma}[Bounded densities preserve the needed stopping integrability]
Let $|\theta_t|\le M$ predictably and let $d\Q=D_Td\Pp$, $D=\cE(\int\theta\,dB)$. If $\sup_\tau\E|X_\tau|^p<\infty$ for one $p>1$, then $X$ is of class (D) under $\Q$.
\end{lemma}
\begin{proof}
Novikov's condition holds for every multiple of $\theta$, and
$\E D_T^q\le\exp(q(q-1)M^2T/2)$ for every $q>1$, by writing $D_T^q$ as the exponential with integrand $q\theta$ times its bounded quadratic-variation correction. Choose $1<r<p$ and put $q=p/(p-r)$. H\"older's inequality gives
\[
 \E_\Q|X_\tau|^r\le
 (\E D_T^q)^{1/q}(\E|X_\tau|^p)^{r/p}.
\]
The right side is uniformly bounded over stopping times. The $r$th-moment bound implies class (D).
\end{proof}

\begin{theorem}[Uniqueness relative to a bounded-integrand solution]
Suppose one admissible solution $(\widetilde Y,\widetilde Z)$ satisfies
$|\widetilde Z_t|\le M$ a.e. for a finite deterministic $M$. Then every other admissible solution with the same terminal value coincides with it.
\end{theorem}
\begin{proof}
For any $|w|\le M$, let $r=|z-w|$. If $r\le1$, the preceding local Lipschitz bound gives
$|g(z)-g(w)|\le(M+2)r$.
If $r>1$, linear growth gives
$|g(z)-g(w)|\le|z|+|w|\le r+2M\le(1+2M)r$.
Thus in both cases
$|g(z)-g(w)|\le(2M+2)|z-w|$.

Set $\Delta Y=Y-\widetilde Y$, $\Delta Z=Z-\widetilde Z$, and define
\[
 b_t=\frac{g(Z_t)-g(\widetilde Z_t)}{|\Delta Z_t|^2}\Delta Z_t
\]
on $\{\Delta Z_t\ne0\}$, zero otherwise. This is predictable and uniformly bounded by $2M+2$. The difference equation under its Girsanov probability is $d\Delta Y=\Delta Z\,dB^\Q$, a local martingale. The power class-(D) condition implies a uniformly bounded stopping $p$th moment, so the previous lemma makes $\Delta Y$ of class (D) under $\Q$. Its terminal value is zero, hence it vanishes. Quadratic variation gives equality of $Z$.
\end{proof}

\begin{corollary}[Explicit affine-terminal subclass]
For $c\in\R$, $z_0\in\R^d$, and $\xi=c+z_0\cdot B_T$, the unique admissible solution is
\[
 Y_t=c+z_0\cdot B_t+(T-t)g(z_0),\qquad Z_t=z_0.
\]
\end{corollary}
\begin{proof}
Substitution verifies the equation. The supremum of this $Y$ has finite moments of every finite order by the Brownian maximal moment bound, or by Doob's inequality at any exponent larger than $p$. Thus $|Y|^p$ is dominated by an integrable variable and is of class (D). The preceding theorem applies to its constant integrand. In particular, constant and zero terminal values cannot produce a counterexample in the required class.
\end{proof}

\begin{theorem}[Bounded-terminal uniqueness]
For every bounded terminal value, the original power class-(D) solution class contains at most one solution.
\end{theorem}
\begin{proof}
Take any solution and define
$\theta_t=g(Z_t)Z_t/|Z_t|^2$ on $\{Z_t\ne0\}$, zero otherwise. Since $|g(Z)|\le|Z|$, one has $|\theta|\le1$ and $g(Z)=\theta\cdot Z$. Under the associated probability $\Q^\theta$, $Y$ is a local martingale. The bounded-density transfer lemma makes it of class (D), so
$Y_t=\E_{\Q^\theta}[\xi\mid\F_t]$.
If $|\xi|\le K$, this proves $|Y_t|\le K$ for every solution, though the probability may depend on the solution.

For every stopping time $\tau$, apply It\^o to $Y^2$ from $\tau$ to localizing times tending to $T$. Using the deterministic bound on $Y$, $|g(Z)|\le|Z|$, and
$2K|Z|\le|Z|^2/2+2K^2$, one obtains
\[
 \E\left[\left.\int_\tau^T|Z_t|^2dt\right|\F_\tau\right]
 \le2K^2+4K^2T.
\]
In detail, before taking the limit the expectation of the integral of $|Z|^2$ is bounded by $K^2+2K\E_\tau\int|Z|$; move half the quadratic term to the left, and then use conditional Fatou. This proves BMO for both solutions' martingale integrands.

For the divided-difference $b$ between two solutions, the local Lipschitz estimate gives
$|b|\le1+|Z|+|\widetilde Z|$, so
$|b|^2\le3(1+|Z|^2+|\widetilde Z|^2)$.
It is therefore BMO. The continuous BMO exponential theorem gives a true density. Under its probability $\Delta Y$ is a bounded local martingale, hence a true martingale closed by zero. This proves equality of $Y$ and then of $Z$.
\end{proof}

\subsection*{5) VERIFICATION}
The explicit sequence gives generator differences $1.021007$, $1.000211$, $1.00000211$, $1.000000021$ for $k=1,10,100,1000$, while its input distances tend to zero. The script additionally samples 30,000 pairs in dimension three with $|w|\le1$; the maximum divided difference is approximately $1.3855$, below the proved bound $4$. Sampling is only a check; the two-case inequality proves the bound for every pair.

The density exponent accommodates $r\in(1,p)$ without assuming a larger terminal moment. Each solution's separate measure only bounds $Y$; comparison needs its own linearization density, obtained for bounded data from the proved BMO estimate.

A cutoff $|Z|,|\widetilde Z|\le R$ produces a Lipschitz constant of order $R$. Usual exponential estimates can then grow like $e^{cR^2T}$, and a mere polynomial terminal moment does not justify discarding the stopping tail against such a factor. No limiting comparison is asserted without an estimate that actually controls this product.

\subsection*{6) FINAL}
\unresolved

\textbf{Strongest checked partial result.} Uniqueness holds for every bounded terminal datum and whenever at least one competing solution has a uniformly bounded integrand. All affine Brownian terminal values, including zero, are explicitly solved in the required class.

\textbf{Exact first gap.} For two arbitrary admissible solutions with unbounded $L^p$ terminal data, justify the change of measure and terminal sampling of their difference for the unbounded divided-difference integrand $b$. Neither BMO nor an adequate density moment was derived here from the fixed power class-(D) hypothesis.

\textbf{Next three moves.} (1) Exploit oscillatory cancellation in localized comparison constants. (2) Prove weighted uniform integrability for stopped difference densities at the given $p$. (3) Construct matching-terminal, unbounded-integrand solutions and verify their full stopping-time power class (D).

\textbf{Likely minimal counterexample, if any.} Both competing integrands and the terminal datum must be unbounded. Large-radius oscillatory excursions would still have to satisfy the full power class-(D) condition.

\clearpage
\begin{thebibliography}{99}
\bibitem{Kazamaki}
Norihiko Kazamaki. \emph{Continuous Exponential Martingales and BMO}. Lecture Notes in Mathematics 1579, Springer, 1994. DOI: 10.1007/BFb0073585. Continuous BMO stochastic exponentials and reverse H\"older inequalities. \url{https://doi.org/10.1007/BFb0073585}.

\bibitem{FHT}
Shengjun Fan, Ying Hu and Shanjian Tang. \emph{1D nonlinear backward stochastic differential equations: a unified theory and applications}. arXiv:2309.06233v2, 11 February 2026. In particular Problems 5.1--5.5, PDF pp.\ 32--33; for C65-45 see (A1), (4.13)--(4.14), PDF p.\ 31. \url{https://arxiv.org/pdf/2309.06233v2}.

\end{thebibliography}
\end{document}
